\chapter{EnvReqs \& Runtime Settings}
\label{ch:envreqs}

\newthought{One block controls the machinery.} \yamlkey{EnvReqs} carries two things: a
V1-heritage entry point for \emph{shared default settings}, and the \yamlkey{V2} map ---
the only place scheduling, \file{SAMPLE/} layout, and archiving policy are configured.
There are deliberately few knobs; this chapter is all of them.

\section{Shared defaults: \yamlkey{Check\_default\_dependencies}}
\label{sec:check-default-deps}

Projects that keep a common environment file can load runtime defaults from it:

\begin{yamlwin}
EnvReqs:
  Check_default_dependencies:
    required: true
    default_yaml_path: "&J/deps/environment_default.yaml"
\end{yamlwin}

When \yamlkey{required} is true, the loader resolves \yamlkey{default\_yaml\_path} (normal
\marker{\&J/} rules), loads that file, and reads \emph{only} its \yamlkey{EnvReqs.V2} map.
Those defaults are recursively merged under the task's own \yamlkey{EnvReqs.V2} --- the
task always wins. With \yamlkey{required: false} the file is skipped; a required but
missing or malformed file fails the load with a focused error.

\begin{jnote}
Other keys in \yamlkey{EnvReqs} --- \yamlkey{OS}, \yamlkey{Python},
\yamlkey{CERN\_ROOT}, and similar dependency declarations --- are preserved as metadata
but not validated at load time. Only \yamlkey{EnvReqs.V2} changes runtime behavior.
\end{jnote}

\section{The \yamlkey{V2} map}
\label{sec:v2-map}

\begin{yamlwin}
EnvReqs:
  V2:
    workers: 4                # Worker processes
    batch_size: 256           # sampler submit-group size
    sample_directory:         # SAMPLE/ bucket policy      (S 5.3)
      enabled: true
      limit: 200
      width: 6
      pack: true
      start_bucket: 1
    cleanup:                  # result handoff policy      (S 5.4)
      strategy: direct        # direct | mv_to_staging
    archiver:                 # persistence policy         (S 5.5)
      mode: process           # process | thread
      handoff: direct
      pack_buckets: true
      batch_size: 200
      flush_interval_sec: 1.0
      delete_after_archive: true
\end{yamlwin}

\begin{jwarn}
\yamlkey{EnvReqs.V2} is validated strictly: an \emph{unknown} key here is a hard error at
load time, so stale settings cannot silently change a run. The \emph{values}, by
contrast, are coerced leniently (Section~\ref{sec:validation}).
\end{jwarn}

\section{Scheduling: \yamlkey{workers} and \yamlkey{batch\_size}}
\label{sec:workers}

\begin{keylist}
  \item[workers] \opt, integer $\geq 0$, default \code{0}. The number of long-lived
        Worker processes. Zero or negative means \emph{one} Worker. The singular alias
        \yamlkey{worker} is accepted; do not write both.
  \item[batch\_size] \opt, integer $> 0$, default \code{256}. How many task
        dictionaries the sampler pushes to Redis per submit group. This is queue
        pipelining only --- execution is always one Sample per Worker --- so the default
        is almost always right.
\end{keylist}

Sizing \yamlkey{workers} is the single most important performance decision you make;
Chapter~\ref{ch:performance} gives the full guidance. The short version: for
calculator-bound scans, set it near the number of calculator invocations your machine (and
license pool) can genuinely sustain.

\section{\yamlkey{sample\_directory}: the SAMPLE bucket policy}
\label{sec:sample-directory}

Controls the bucketed layout introduced in Section~\ref{sec:sample-lifecycle}:
\file{SAMPLE/\allowbreak{}<bucket>/\allowbreak{}<uuid>/}, packed to \file{<bucket>.\allowbreak{}tar.\allowbreak{}gz} after archiving. The
same map is also accepted as \yamlkey{Scan.\allowbreak{}sample\_directory}.

\begin{table}[ht]
  \footnotesize
  \begingroup
  \setlength{\tabcolsep}{0pt}
  \begin{tabular}{@{}p{0.24\linewidth}p{0.20\linewidth}p{0.50\linewidth}@{}}
    \toprule
    \textbf{Key} & \textbf{Default} & \textbf{Meaning} \\
    \midrule
    \yamlkey{enabled} & \code{true} & allocate Samples to buckets on pull \\
    \yamlkey{limit}   & \code{200}  & Samples per bucket directory \\
    \yamlkey{width}   & \code{6}    & zero-pad width of bucket names (\file{000001}) \\
    \yamlkey{pack}    & \code{true} & tar a bucket once fully archived \\
    \yamlkey{start\_bucket} & \code{1} & first bucket number \\
    \bottomrule
  \end{tabular}
  \endgroup
  \caption{\yamlkey{sample\_directory} keys.}
  \label{tab:sample-directory}
\end{table}

A bucket is packed only when it is sealed, no Sample in it is still active, and every
Sample in it has been archived to the database --- packing can never outrun persistence.

\section{\yamlkey{cleanup}: result handoff}
\label{sec:cleanup}

\begin{keylist}
  \item[strategy] \code{direct} (default) or \code{mv\_to\_staging}. With
        \code{direct}, a finished Sample's products are already in their final home
        under \file{SAMPLE/} and the Archiver only writes database records. With
        \code{mv\_to\_staging}, the Worker first moves the Sample directory into a
        staging area and the Archiver moves it onward --- a buffer hop for setups where
        final storage is slow or remote.
\end{keylist}

The same map may be written as \yamlkey{Calculators.Cleanup}, where a
\yamlkey{staging\_dir} override is also available (default
\file{<task\_result\_dir>/staging}); the task-card value wins over the
\yamlkey{EnvReqs.V2} value.

\section{\yamlkey{archiver}: persistence policy}
\label{sec:archiver-policy}

\begin{table}[ht]
  \footnotesize
  \begingroup
  \setlength{\tabcolsep}{0pt}
  \begin{tabular}{@{}p{0.24\linewidth}p{0.20\linewidth}p{0.50\linewidth}@{}}
    \toprule
    \textbf{Key} & \textbf{Default} & \textbf{Meaning} \\
    \midrule
    \yamlkey{mode} & \code{process} & dedicated Archiver process, or a thread in the
                                       control process (\code{thread}) \\
    \yamlkey{handoff} & \code{direct} & mirrors \yamlkey{cleanup.strategy} \\
    \yamlkey{pack\_buckets} & \code{true} & tar sealed buckets after archiving \\
    \yamlkey{batch\_size} & \code{200} & Samples buffered per database flush \\
    \yamlkey{flush\_interval\_sec} & \code{1.0} & time-based flush for partial batches
                                                  (floor \code{0.05}) \\
    \yamlkey{strategy} & \code{move} & \code{move} or \code{copy} when a staging hop is
                                       active \\
    \yamlkey{delete\_after\_archive} & \code{true} & remove staging residue when a
                                                     staging hop is active \\
    \bottomrule
  \end{tabular}
  \endgroup
  \caption{\yamlkey{archiver} keys. Also writable as \yamlkey{Calculators.Archiver};
  task values win.}
  \label{tab:archiver-keys}
\end{table}

The defaults --- a dedicated process, direct handoff, batched flushes, packed buckets ---
are the production configuration; Chapter~\ref{ch:distributed-runtime} explains the
machinery behind them and when to deviate.

\section{What is \emph{not} here}
\label{sec:not-here}

Deliberate omissions, so you do not go looking for them:

\begin{itemize}
  \item \textbf{Redis connection settings.} The broker is always the local service on
        the default port (Chapter~\ref{ch:installation}).
  \item \textbf{Per-sampler thread counts or worker pinning.} Workers compete on one
        shared queue; there is nothing to pin.
  \item \textbf{A mode switch.} There is exactly one runtime path
        (Section~\ref{sec:runtime-picture}).
\end{itemize}
